%%% TITLE, AUTHOR, AFFILIATIONS, ABSTRACT %%%

%%% ------------ TITLE ------------ %%%

\title{Methylation Clocks Fail to Generalize Across Genetically Admixed Individuals}

%%% ------------ AUTHORS ------------ %%%

\author[1,2]{Sebastián~Cruz-González}
\author[3]{Ogechukwu Okpala}
\author[4]{Esther~Gu}
\author[4]{Lissette~Gomez}
\author[5]{Makaela~Mews}
\author[4,6]{Jeffery~M.~Vance}
\author[4,6]{Michael~L.~Cuccaro}
\author[7]{Mario~R.~Cornejo-Olivas}
\author[8]{Briseida~E.~Feliciano-Astacio}
\author[9]{Goldie~S.~Byrd}
\author[5]{Jonathan~L.~Haines}
\author[4,6]{Margaret~A.~Pericak-Vance}
\author[4]{Anthony~J.~Griswold}
\author[5]{William~S.~Bush}
\author[2,3,10,*]{John~A.~Capra}
\affil[1]{Biological and Medical Informatics Program, University of California, San Francisco, San Francisco, CA}
\affil[2]{Bakar Computational Health Sciences Institute, University of California, San Francisco, CA}
\affil[3]{Department of Epidemiology and Biostatistics, University of California, San Francisco, CA}
\affil[4]{John P. Hussman Institute for Human Genomics, Miller School of Medicine, University of Miami, Miami, FL}
\affil[5]{Department of Population and Quantitative Health Sciences, Case Western Reserve University, Cleveland, OH}
\affil[6]{The Dr. John T. Macdonald Foundation Department of Human Genetics, Miller School of Medicine, University of Miami, Miami, FL}
\affil[7]{Neurogenetics Research Center, Instituto Nacional de Ciencias Neurologicas, Lima, 15003, Peru}
\affil[8]{Department of Internal Medicine, Universidad Central Del Caribe, Bayamón, Puerto Rico}
\affil[9]{Maya Angelou Center for Health Equity, Wake Forest University, Winston-Salem, NC}
\affil[10]{Department of Bioengineering and Therapeutic Sciences, University of California, San Francisco, CA}

\affil[ ]{ } % for some blank space
\affil[*]{\textit{Correspondence to }\underline{tony@capralab.org}}
\renewcommand\Affilfont{\footnotesize} % size of affilitations
\setcounter{Maxaffil}{0} % no max for num affiliations
\date{} % If you want a date fill this in (e.g. \date{February 2022})

%%% ------------ ABSTRACT ------------ %%%

\newcommand{\makeAbstract}{
\begin{abstract}
\noindent
Epigenetic aging clocks based on DNA methylation patterns across the genome have emerged as a potential biomarker for risk of age-related diseases, like Alzheimer’s disease (AD), and environmental and social stressors. However, methylation clocks have not been comprehensively validated in genetically diverse individuals. Here we evaluate a set of first-, second-, and third-generation methylation clocks in 621 AD patients and matched controls from African American, Hispanic, and White cohorts. The clocks are less accurate at predicting age in genetically admixed cohorts compared to the White cohort, especially for those with substantial African ancestry. This decreased accuracy holds in $>$2,500 individuals of European and African ancestry from three additional datasets. The clocks also fail to consistently identify age acceleration in admixed AD cases compared to controls. To explore potential causes for the lack of generalization of the clocks, we intersected clock CpGs with methylation, germline genetic variants, and methylation QTL (meQTL) data from global populations. We find differential methylation between African and European ancestry individuals is common for clock CpGs. Genetic variants rarely disrupt clock CpGs between populations, but a substantial fraction of clock CpGs have meQTL with significantly higher frequencies in African genetic ancestries. Our results demonstrate that methylation clocks often fail to predict age and AD risk when applied across populations and suggest avenues for improving their portability by considering differences in genetic and epigenetic patterns across human populations.
\end{abstract}
}

\clearpage
